\documentclass[10pt,a4paper]{amsart}
\usepackage{amsmath,amssymb,amsthm,graphicx}
% Retained preamble for the Stacks Project chapter {CH} (shared project preamble).
% The chapter file pinned-original-source.tex only carries \input{{preamble}} on its first line,
% so no preamble text lives inside the chapter itself.  The Stacks Project uses one single shared
% preamble for every chapter: the snapshot stages it three times (sources/stacks-chow/preamble.tex,
% sources/stacks-etale-cohomology/preamble.tex and focused-seeds/AGP/sources/stacks-snapshot-20261011/
% preamble.tex) and those copies are identical in size.  The block below is copied mechanically,
% line for line, from sources/stacks-chow/preamble.tex lines 14-18 and 154-217: the Xy-pic setup and
% the native theorem declarations together with the mathematical macros.  Nothing was added, removed
% or redefined; the wrapper class and the amsmath/amssymb/amsthm/graphicx base packages are supplied
% by the assembler.
% For commutative diagrams we use Xy-pic
% Chapter: stacks-spaces-more-morphisms.tex -- More on Morphisms of Spaces.
\usepackage[all]{xy}

% We use 2cell for 2-commutative diagrams.
\xyoption{2cell}

% Theorem environments.
%
\theoremstyle{plain}
\newtheorem{theorem}[subsection]{Theorem}
\newtheorem{proposition}[subsection]{Proposition}
\newtheorem{lemma}[subsection]{Lemma}

\theoremstyle{definition}
\newtheorem{definition}[subsection]{Definition}
\newtheorem{example}[subsection]{Example}
\newtheorem{exercise}[subsection]{Exercise}
\newtheorem{situation}[subsection]{Situation}

\theoremstyle{remark}
\newtheorem{remark}[subsection]{Remark}
\newtheorem{remarks}[subsection]{Remarks}

\numberwithin{equation}{subsection}

% Macros
%
\def\lim{\mathop{\mathrm{lim}}\nolimits}
\def\colim{\mathop{\mathrm{colim}}\nolimits}
\def\Spec{\mathop{\mathrm{Spec}}}
\def\Hom{\mathop{\mathrm{Hom}}\nolimits}
\def\Ext{\mathop{\mathrm{Ext}}\nolimits}
\def\SheafHom{\mathop{\mathcal{H}\!\mathit{om}}\nolimits}
\def\SheafExt{\mathop{\mathcal{E}\!\mathit{xt}}\nolimits}
\def\Sch{\mathit{Sch}}
\def\Mor{\mathop{\mathrm{Mor}}\nolimits}
\def\Ob{\mathop{\mathrm{Ob}}\nolimits}
\def\Sh{\mathop{\mathit{Sh}}\nolimits}
\def\NL{\mathop{N\!L}\nolimits}
\def\CH{\mathop{\mathrm{CH}}\nolimits}
\def\proetale{{pro\text{-}\acute{e}tale}}
\def\etale{{\acute{e}tale}}
\def\QCoh{\mathit{QCoh}}
\def\Ker{\mathop{\mathrm{Ker}}}
\def\Im{\mathop{\mathrm{Im}}}
\def\Coker{\mathop{\mathrm{Coker}}}
\def\Coim{\mathop{\mathrm{Coim}}}

% Boxtimes
%
\DeclareMathSymbol{\boxtimes}{\mathbin}{AMSa}{"02}

%
% Macros for moduli stacks/spaces
%
\def\QCohstack{\mathcal{QC}\!\mathit{oh}}
\def\Cohstack{\mathcal{C}\!\mathit{oh}}
\def\Spacesstack{\mathcal{S}\!\mathit{paces}}
\def\Quotfunctor{\mathrm{Quot}}
\def\Hilbfunctor{\mathrm{Hilb}}
\def\Curvesstack{\mathcal{C}\!\mathit{urves}}
\def\Polarizedstack{\mathcal{P}\!\mathit{olarized}}
\def\Complexesstack{\mathcal{C}\!\mathit{omplexes}}
% \Pic is the operator that assigns to X its picard group, usage \Pic(X)
% \Picardstack_{X/B} denotes the Picard stack of X over B
% \Picardfunctor_{X/B} denotes the Picard functor of X over B
\def\Pic{\mathop{\mathrm{Pic}}\nolimits}
\def\Picardstack{\mathcal{P}\!\mathit{ic}}
\def\Picardfunctor{\mathrm{Pic}}
\def\Deformationcategory{\mathcal{D}\!\mathit{ef}}
\usepackage{cleveref}
\crefname{originallemma}{original lemma}{original lemmas}
\Crefname{originallemma}{Original lemma}{Original lemmas}
\makeatletter
\newlabel{lemma-flat-after-blowing-up}{{087E}{source}{}{sourcepointer-1}{}}
\newlabel{lemma-flat-after-blowing-up@cref}{{[originallemma][0][]087E}{[1][0][]source}{}{}{}}
\makeatother
\crefname{originallemma}{original lemma}{original lemmas}
\Crefname{originallemma}{Original lemma}{Original lemmas}
\makeatletter
\newlabel{lemma-proper-flat-finite-over-dense-open}{{0B4J}{source}{}{sourcepointer-2}{}}
\newlabel{lemma-proper-flat-finite-over-dense-open@cref}{{[originallemma][0][]0B4J}{[1][0][]source}{}{}{}}
\makeatother
\crefname{originallemma}{original lemma}{original lemmas}
\Crefname{originallemma}{Original lemma}{Original lemmas}
\makeatletter
\newlabel{lemma-zariski}{{0879}{source}{}{sourcepointer-3}{}}
\newlabel{lemma-zariski@cref}{{[originallemma][0][]0879}{[1][0][]source}{}{}{}}
\makeatother
\title{A proper morphism that is an isomorphism over a quasi-compact open becomes an isomorphism after an admissible blowup}
\author{The Stacks Project Authors}
\date{stacks-project:spaces-more-morphisms.tex:local-snapshot-20261011; GNU FDL 1.2 or later}
\begin{document}
\maketitle
\noindent\textit{Editorial note.} Counted span is the author's own proof environment, source lines 8950-8959, immediately following the statement at 8936-8948.\par\medskip
\noindent\textit{Editorial note.} This chapter (spaces-more-morphisms) has no entry in reference/public-items.jsonl, so the range check is vacuous for it; the cross-lane label scan is the substantive check.\par\medskip
\noindent\textit{Editorial note.} Three chapter lemmas are cited by \textbackslash{}ref as navigation pointers into the same chapter.\par\medskip
\noindent\textit{Allowed standard prerequisites (editorial).}\begin{itemize}
\item Flatness and finite presentation after an admissible blowup. chapter lemma lemma-flat-after-blowing-up, cited at proof line 8951. Allows replacing the morphism by one that is flat and of finite presentation, at the cost of a U-admissible blowup.
\item A proper, flat morphism over a scheme-theoretically dense open is finite. chapter lemma lemma-proper-flat-finite-over-dense-open, cited at proof line 8955. Turns properness plus flatness plus density of U into finiteness of f.
\item Zariski's main theorem style criterion. chapter lemma lemma-zariski, cited at proof line 8956. Turns finiteness plus the fact that f is an isomorphism over U into an open immersion.
\item Existence of a U-admissible blowup making U scheme-theoretically dense. used at proof lines 8953-8954. Produces the blowup after which the density hypothesis holds.
\end{itemize}\medskip
\ifdefined\hypertarget\hypertarget{sourcepointer-1}{}\fi
\noindent\textit{Source reference.} Original lemma 087E: sources/stacks-spaces-more-morphisms/pinned-original-source.tex, source lines 8861--8861. Navigation pointer only: the counted proof cites this statement by \textbackslash{}ref as a standing input; the deduction the counted proof has to make is carried out inside the counted span, and no part of that deduction is delegated to this pointer.\par\medskip
\ifdefined\hypertarget\hypertarget{sourcepointer-2}{}\fi
\noindent\textit{Source reference.} Original lemma 0B4J: sources/stacks-spaces-more-morphisms/pinned-original-source.tex, source lines 8625--8625. Navigation pointer only: the counted proof cites this statement by \textbackslash{}ref as a standing input; the deduction the counted proof has to make is carried out inside the counted span, and no part of that deduction is delegated to this pointer.\par\medskip
\ifdefined\hypertarget\hypertarget{sourcepointer-3}{}\fi
\noindent\textit{Source reference.} Original lemma 0879: sources/stacks-spaces-more-morphisms/pinned-original-source.tex, source lines 8649--8649. Navigation pointer only: the counted proof cites this statement by \textbackslash{}ref as a standing input; the deduction the counted proof has to make is carried out inside the counted span, and no part of that deduction is delegated to this pointer.\par\medskip
\begin{lemma}
\label{lemma-isomorphism-after-blowing-up}
Let $S$ be a scheme. Let $f : X \to B$ be a morphism of algebraic
spaces over $S$. Let $U \subset B$ be an open subspace. Assume
\begin{enumerate}
\item $B$ is quasi-compact and quasi-separated,
\item $U$ is quasi-compact,
\item $f : X \to B$ is proper, and
\item $f^{-1}(U) \to U$ is an isomorphism.
\end{enumerate}
Then there exists a $U$-admissible blowup $B' \to B$ such that
the strict transform $X'$ of $X$ maps isomorphically to $B'$.
\end{lemma}

\begin{proof}
By Lemma \ref{lemma-flat-after-blowing-up} we may assume that
$X \to B$ is flat and of finite presentation. After replacing
$B$ by a $U$-admissible blowup if necessary, we may assume
that $U \subset B$ is scheme theoretically dense. Then $f$ is
finite by Lemma \ref{lemma-proper-flat-finite-over-dense-open}
and an open immersion by Lemma \ref{lemma-zariski}. Hence $f$
is an open immersion whose image is closed and contains the
dense open $U$, whence $f$ is an isomorphism.
\end{proof}

\end{document}
